% =====================================================================
% Group E -- Possible further questions: topics marked T / T&N in the course program that have not been asked so far
% Same format, shorter answers.
% =====================================================================
\qgroup{E}{Possible further questions (T topics never asked so far)}

Topics marked T or T\&N in \texttt{Exam\_Course\_INFO.pdf} that do not appear in the collected exams. They are grouped by
exam slot: E1--E13 can be the "third" orbital question (chapters 2, 3, 4, 5, 6), E14 a Chapter 8 question, E15--E19 the
attitude question. Answers are shorter than in groups A--D but follow the same structure.

\begin{center}
\small
\begin{tabular}{@{}llr@{\hspace{1.2em}}llr@{}}
\toprule
Q & Topic & Page & Q & Topic & Page\\
\midrule
\ref{q:E1} & Two-body problem, potential of a sphere & \pageref{q:E1} & \ref{q:E11} & Newton's law with propulsion & \pageref{q:E11}\\
\ref{q:E2} & Eccentricity vector & \pageref{q:E2} & \ref{q:E12} & Optimal staging & \pageref{q:E12}\\
\ref{q:E3} & Energy, vis-viva, cosmic velocities & \pageref{q:E3} & \ref{q:E13} & Ascent trajectory & \pageref{q:E13}\\
\ref{q:E4} & Kepler's equation & \pageref{q:E4} & \ref{q:E14} & Perturbations: overview & \pageref{q:E14}\\
\ref{q:E5} & Special trajectories & \pageref{q:E5} & \ref{q:E15} & Euler's principal rotation theorem & \pageref{q:E15}\\
\ref{q:E6} & ECI, LVLH, orbit elements & \pageref{q:E6} & \ref{q:E16} & Euler's equations, kinetic energy & \pageref{q:E16}\\
\ref{q:E7} & Ground track properties & \pageref{q:E7} & \ref{q:E17} & Inertia matrix: properties, symmetry & \pageref{q:E17}\\
\ref{q:E8} & Three-dimensional transfers & \pageref{q:E8} & \ref{q:E18} & Torque-free axisymmetric body & \pageref{q:E18}\\
\ref{q:E9} & Sphere of influence, flyby & \pageref{q:E9} & \ref{q:E19} & Manoeuvres of spinning satellites & \pageref{q:E19}\\
\ref{q:E10} & Problem of two bodies & \pageref{q:E10} & & & \\
\bottomrule
\end{tabular}
\end{center}

% ---------------------------------------------------------------------
\question{E1}{Two-body problem and potential of a spherical mass distribution}
% source: cap02.tex l.35-166
\begin{qbox}
State the fundamental principles of orbital mechanics, derive the equation of the restricted two-body problem and show that a
spherical mass distribution generates the same potential as a point mass.
\qmeta{never asked (Ch.\,2, T).}
\end{qbox}
\begin{outline}
\begin{enumerate}
\item Newton's 2nd law (inertial frame) and universal gravitation; equivalence principle.
\item Two-body problem: subtract the equations $\Rightarrow$ $\frac{d^2}{dt^2}(\underline R_2-\underline R_1)=-\frac{G(m_1+m_2)}{R_{12}^3}(\underline R_2-\underline R_1)$.
\item Restricted ($m_1\gg m_2$): $\ddot{\underline r}=-\frac{\mu}{r^3}\underline r$, $\mu=Gm_1$.
\item Potential $V=\mu/r$ ($\underline F=\nabla V$, $U=-V$); thin spherical shell $\Rightarrow$ same as a point mass at the centre; sum of shells.
\end{enumerate}
\end{outline}
\answer
\paragraph{Principles.} (A) Newton's 2nd law in an inertial frame, $\underline F=\frac{d}{dt}(m\underline v)$. (B) Universal gravitation:
$\underline F_1=G\frac{m_1m_2}{R_{12}^3}(\underline R_2-\underline R_1)=-\underline F_2$, $G=6.673\cdot10^{-11}$\,m$^3$kg$^{-1}$s$^{-2}$. The equivalence
principle allows identifying inertial and gravitational mass.

\paragraph{Two-body problem.} Writing $m_1\ddot{\underline R}_1=\underline F_1$, $m_2\ddot{\underline R}_2=\underline F_2$ and subtracting the accelerations:
\begin{equation}
\frac{d^{2}}{dt^{2}}\left(\underline{R}_{2}-\underline{R}_{1}\right)=-G\frac{m_{1}+m_{2}}{R_{12}^{3}}\left(\underline{R}_{2}-\underline{R}_{1}\right),
\end{equation}
the motion of body 2 relative to body 1. If $m_1\gg m_2$ (\textbf{restricted} two-body problem), $m_1+m_2\simeq m_1$ and the effect of 2 on
1 is negligible; with $\underline r=\underline R_2-\underline R_1$:
\begin{equation}
\frac{d^{2}\underline r}{dt^{2}}=-\frac{\mu}{r^{3}}\underline r,\qquad \mu=Gm_1\ (\mu_\oplus=398600.4\ \text{km}^3/\text{s}^2).
\end{equation}

\paragraph{Gravitational potential.} A conservative force derives from a potential $V$ ($\underline F=\nabla V$) or a potential energy $U=-V$
($\underline F=-\nabla U$). For a point mass, per unit mass, $V=\mu/r$, $U=-\mu/r$.

\paragraph{Spherical mass distribution.} A thin shell of radius $R$, uniform surface density $\sigma=m_1/(4\pi R^2)$; the ring at angle $\theta$
has $dm=2\pi R^2s_\theta d\theta\,\sigma$ at distance $\rho$ from the external point, $\rho^2=r^2+R^2-2rRc_\theta$. Then
$dV=Gm_2\,dm/\rho$ and, using $2\rho\,d\rho=2rRs_\theta d\theta$,
\begin{equation}
V=Gm_2\,2\pi R^2\sigma\int_{r-R}^{r+R}\frac{d\rho}{rR}=\frac{4\pi R^2\sigma\,Gm_2}{r}=G\frac{m_1m_2}{r}.
\end{equation}
A shell generates the same potential as a point mass at its centre; a spherical body is a sum of shells, so it does the same. This
justifies treating spherical planets as point masses (and is the starting point of the harmonic expansion when the body is not spherical).

\begin{keyf}
$\ddot{\underline r}=-\dfrac{\mu}{r^3}\underline r$, \ $\mu=G(m_1+m_2)\simeq Gm_1$, \ $V=\dfrac{\mu}{r}$, $U=-\dfrac{\mu}{r}$, \ shell: $V=G\dfrac{m_1m_2}{r}$ for $r>R$.
\end{keyf}
\pitfalls
\begin{itemize}
\item The relative motion uses $G(m_1+m_2)$; the restricted approximation replaces it with $Gm_1$.
\item The shell result holds for external points ($r>R$).
\end{itemize}
\source{cap02.tex l.35--95 (principles, two-body), l.96--166 (potential, spherical shells). Formulario: --.}

% ---------------------------------------------------------------------
\question{E2}{The eccentricity vector}
% source: cap02.tex l.203-251
\begin{qbox}
Prove that the eccentricity vector is a first integral of the restricted two-body problem and use it to obtain the polar equation of the orbit.
\qmeta{never asked (Ch.\,2, T).}
\end{qbox}
\begin{outline}
\begin{enumerate}
\item $\frac{d}{dt}(\underline h\times\underline v)=-\mu\,\frac{d\hat r}{dt}$.
\item $\underline e:=-\hat r+\frac{\underline v\times\underline h}{\mu}$ $\Rightarrow$ $\dot{\underline e}=\underline 0$.
\item $\underline h\cdot\underline e=0$: 5 independent integrals; $\underline e$ in the plane, towards periapsis.
\item $\underline r\cdot\underline e$ $\Rightarrow$ $r=\frac{h^2/\mu}{1+e\cos\theta_*}$.
\end{enumerate}
\end{outline}
\answer
\paragraph{Proof.} With $\underline h$ constant,
$\frac{d}{dt}(\underline h\times\underline v)=\underline h\times\dot{\underline v}=-\frac{\mu}{r^2}(\underline r\times\underline v)\times\hat r$.
With $(\underline a\times\underline b)\times\underline c=\underline b(\underline a\cdot\underline c)-\underline a(\underline b\cdot\underline c)$:
$(\underline r\times\underline v)\times\hat r=r\underline v-\dot r\,\underline r=r^2\frac{d\hat r}{dt}$ (since $\frac{d\hat r}{dt}=\frac{\underline v-\dot r\hat r}{r}$). Hence
\begin{equation}
\frac{d}{dt}\left(\underline h\times\underline v\right)=-\mu\frac{d\hat r}{dt}
\;\Longrightarrow\;
\underline{e}:=-\hat{r}+\frac{\underline{v}\times\underline{h}}{\mu},\qquad \frac{d\underline e}{dt}=-\frac{d\hat r}{dt}+\frac{d\hat r}{dt}=\underline 0 .
\end{equation}
(The notes reach the same result writing $\underline v=\dot r\hat r+\underline\omega\times\underline r$.)

\paragraph{Properties.} $\underline h\cdot\underline e=0$: $\underline e$ lies in the orbit plane and is inertially fixed. $\underline h$ and $\underline e$ give 5 (not 6)
independent scalar integrals. At periapsis ($\underline v\perp\underline r$), $\underline v\times\underline h=v_ph\,\hat r$, so $\underline e=(v_ph/\mu-1)\hat r_p$: $\underline e$
points towards the periapsis and $|\underline e|=e$ is the eccentricity.

\paragraph{Polar equation.} Let $\theta_*$ (true anomaly) be the angle from $\underline e$ to $\underline r$; for $h\neq0$:
\begin{equation}
\underline r\cdot\underline e=re\cos\theta_*=-r+\frac{\underline r\cdot(\underline v\times\underline h)}{\mu}=-r+\frac{h^2}{\mu}
\;\Longrightarrow\; r=\frac{p}{1+e\cos\theta_*},\quad p=\frac{h^2}{\mu},
\end{equation}
($\underline r\cdot(\underline v\times\underline h)=\underline h\cdot(\underline r\times\underline v)=h^2$): a conic with focus at the attracting body.
Also $\underline h\times\underline e$ gives the velocity $\underline v=\sqrt{\mu/p}\,(e\,\hat p+\hat\theta)$.
\begin{keyf}
$\underline e=-\hat r+\dfrac{\underline v\times\underline h}{\mu}$, \ $\dot{\underline e}=\underline 0$, \ $\underline e\cdot\underline h=0$, \ $r=\dfrac{h^2/\mu}{1+e\,c_{\theta_*}}$, \ $\underline v=\sqrt{\mu/p}\,(e\hat p+\hat\theta)$.
\end{keyf}
\pitfalls
\begin{itemize}
\item Sign convention: $-\hat r+\underline v\times\underline h/\mu$ (order $\underline v\times\underline h$).
\item Say why there are only 5 independent integrals.
\end{itemize}
\source{cap02.tex l.203--228 (eccentricity vector), l.230--251 (polar equation), l.372--410 (velocity). Formulario: \fml{2-HEV}, \fml{2-POL}.}

% ---------------------------------------------------------------------
\question{E3}{Specific energy, vis-viva equation and cosmic velocities}
% source: cap02.tex l.372-587, 929-945
\begin{qbox}
Define the specific energy of a Keplerian orbit, derive its expression in terms of the orbit elements, and define the cosmic velocities.
\qmeta{never asked (Ch.\,2, T\&N).}
\end{qbox}
\begin{outline}
\begin{enumerate}
\item Velocity components $v_r=\sqrt{\mu/p}\,e\,s_{\theta_*}$, $v_\theta=\sqrt{\mu/p}(1+ec_{\theta_*})$.
\item $\mathcal E=\frac{v^2}{2}-\frac{\mu}{r}=-\frac{\mu}{2p}(1-e^2)=-\frac{\mu}{2a}$; sign vs orbit type.
\item Vis-viva $a=(2/r-v^2/\mu)^{-1}$; $e=\sqrt{1-h^2/(\mu a)}$.
\item Cosmic velocities $v_I=\sqrt{\mu/r}$, $v_{II}=\sqrt{2\mu/r}$; $v_\infty$ for hyperbolas.
\end{enumerate}
\end{outline}
\answer
\paragraph{Energy.} $\mathcal E=\frac{v^2}{2}-\frac{\mu}{r}$ (kinetic + potential, per unit mass) is constant (gravity is conservative). With
$r=\frac{p}{1+ec_{\theta_*}}$ and $v^2=v_r^2+v_\theta^2=\frac{\mu}{p}(1+e^2+2ec_{\theta_*})$:
\begin{equation}
\mathcal E=\frac{\mu}{2p}\left(1+e^2+2ec_{\theta_*}\right)-\frac{\mu}{p}\left(1+ec_{\theta_*}\right)=-\frac{\mu}{2p}\left(1-e^2\right)=-\frac{\mu}{2a}.
\end{equation}
$\mathcal E<0$ ellipse ($|PE|>KE$), $=0$ parabola, $>0$ hyperbola ($|PE|<KE$); it depends only on $a$.

\paragraph{Vis-viva.} From $\frac{v^2}{2}-\frac{\mu}{r}=-\frac{\mu}{2a}$: $v^2=\mu\left(\frac2r-\frac1a\right)$, $a=\frac{\mu}{2\mu/r-v^2}$; with
$h=|\underline r\times\underline v|=\sqrt{\mu a(1-e^2)}$: $e=\sqrt{1-h^2/(\mu a)}$. Hence $r$ and $v$ at one point fix $a$ (size, energy) and $e$ (shape).
Special values: ellipse $v_{p,A}=\sqrt{\frac{\mu}{a}\frac{1\pm e}{1\mp e}}$; hyperbola $v_\infty=\sqrt{-\mu/a}=\sqrt{2\mathcal E}$ (hyperbolic excess
velocity); parabola $v=\sqrt{2\mu/r}\to0$.

\paragraph{Cosmic velocities} at a point at distance $r$:
\begin{itemize}
\item \textbf{first}, $v_I=\sqrt{\mu/r}$: horizontal velocity giving a circular orbit of radius $r$ ($\mathcal E=-\mu/2r$);
\item \textbf{second} (escape), $v_{II}=\sqrt{2\mu/r}=\sqrt2\,v_I$: gives $\mathcal E=0$, a parabola with periapsis $r$: escape with zero velocity at infinity.
\end{itemize}
At the Earth's surface $v_I\simeq7.9$\,km/s, $v_{II}\simeq11.2$\,km/s \notinnotes. See \qref{A3} for the full family of trajectories.
\begin{keyf}
$\mathcal E=\dfrac{v^2}{2}-\dfrac{\mu}{r}=-\dfrac{\mu}{2a}$, \ $v^2=\mu\left(\dfrac2r-\dfrac1a\right)$, \ $v_I=\sqrt{\dfrac{\mu}{r}}$, $v_{II}=\sqrt{\dfrac{2\mu}{r}}$, \ $v_\infty=\sqrt{-\dfrac{\mu}{a}}$.
\end{keyf}
\pitfalls
\begin{itemize}
\item $\mathcal E$ does not depend on $e$: orbits with the same $a$ have the same energy and period.
\item $a<0$ for hyperbolas, so that $\mathcal E=-\mu/2a>0$.
\end{itemize}
\source{cap02.tex l.372--530 (velocity), l.532--552 (cosmic velocities), l.554--587 (energy), l.929--945 (vis-viva). Formulario: \fml{2-ENE}, \fml{2-VIS}, \fml{2-COS}.}

% ---------------------------------------------------------------------
\question{E4}{Position in time: eccentric anomaly and Kepler's equation}
% source: cap02.tex l.589-792
\begin{qbox}
Introduce the eccentric anomaly and derive Kepler's equation for elliptic orbits; explain how it is solved.
\qmeta{never asked (Ch.\,2, T\&N).}
\end{qbox}
\begin{outline}
\begin{enumerate}
\item $\dot\theta_*=h/r^2$: integrable only numerically $\Rightarrow$ alternative via $E$.
\item Definition of $E$ (auxiliary circle): $x=ac_E$, $y=bs_E$; $\tan\frac E2=\sqrt{\frac{1-e}{1+e}}\tan\frac{\theta_*}{2}$; $r=a(1-ec_E)$.
\item $\dot E(1-ec_E)=\sqrt{\mu/a^3}$ $\Rightarrow$ $M=E-es_E$, $M-M_0=n(t-t_0)$.
\item Newton's method; period $T=2\pi\sqrt{a^3/\mu}$ (3rd Kepler law).
\end{enumerate}
\end{outline}
\answer
\paragraph{Motivation.} $\dot\theta_*=\frac{h}{r^2}=\sqrt{\frac{\mu}{p^3}}(1+ec_{\theta_*})^2$ relates true anomaly and time but must be integrated
numerically; for ellipses a cheaper method uses the \textbf{eccentric anomaly} $E$.

\paragraph{Eccentric anomaly.} In a frame centred at the centre of the ellipse ($x$ along the apse line), the spacecraft is at
$x=ea+rc_{\theta_*}$, $y=rs_{\theta_*}$; $E$ is defined by $x=ac_E$, $y=bs_E=a\sqrt{1-e^2}\,s_E$ (projection on the auxiliary circle). Then
\begin{equation}
\tan\frac E2=\frac{s_E}{1+c_E}=\sqrt{\frac{1-e}{1+e}}\tan\frac{\theta_*}{2},\qquad r=a\left(1-e\,c_E\right).
\end{equation}

\paragraph{Kepler's equation.} Differentiating the two expressions of $r$: $\dot r=ae\dot Es_E$ and $\dot r=\sqrt{\mu/a^3}\,\frac{ae\,s_{\theta_*}}{\sqrt{1-e^2}}$;
with $rs_{\theta_*}=bs_E$ this gives $\dot E(1-ec_E)=\sqrt{\mu/a^3}$, which integrates (separation of variables) to
\begin{equation}
\left(E-e\,s_{E}\right)-\left(E_{0}-e\,s_{E_{0}}\right)=\sqrt{\frac{\mu}{a^{3}}}\,(t-t_{0})
\quad\Longleftrightarrow\quad M-M_0=n\,(t-t_0),\qquad M:=E-e\,s_E,\ n=\sqrt{\mu/a^3}.
\end{equation}
The \textbf{mean anomaly} $M$ grows linearly in time. Time from position: $\theta_*\to E\to M\to t$ (explicit). Position from time:
$t\to M\to E$ (transcendental, numerical) $\to\theta_*$. Circular orbits: $\theta_*=E=M$.

\paragraph{Newton's method.} $f(E)=E-es_E-M=0$, $E_{k+1}=E_k-\frac{f(E_k)}{1-ec_{E_k}}$, starting from $E_1=M$; converges in a few iterations.
After one period $\Delta E=2\pi=nT$: $T=2\pi\sqrt{a^3/\mu}$, i.e.\ $T^2/a^3=4\pi^2/\mu$ (3rd Kepler law). Analogous laws: Barker's equation for
parabolas, hyperbolic anomaly for hyperbolas.
\begin{keyf}
$\tan\frac E2=\sqrt{\frac{1-e}{1+e}}\tan\frac{\theta_*}2$, \ $r=a(1-ec_E)$, \ $M=E-es_E=M_0+\sqrt{\mu/a^3}(t-t_0)$, \ $E_{k+1}=E_k-\dfrac{E_k-es_{E_k}-M}{1-ec_{E_k}}$, \ $T=2\pi\sqrt{a^3/\mu}$.
\end{keyf}
\pitfalls
\begin{itemize}
\item Use the half-angle formula to keep $E$ and $\theta_*$ in the same half-plane.
\item $E$, $M$ in radians; count full revolutions separately.
\end{itemize}
\source{cap02.tex l.589--615 ($\dot\theta_*$), l.617--748 (eccentric anomaly, Kepler's equation, Newton's method, 3rd law). Formulario: \fml{2-E2T}, \fml{2-KEP}, \fml{2-NEW}.}

% ---------------------------------------------------------------------
\question{E5}{Special Keplerian trajectories}
% source: cap02.tex l.794-891
\begin{qbox}
Describe the special Keplerian trajectories: circular orbits (with the balance of accelerations), geosynchronous, ballistic and rectilinear trajectories.
\qmeta{never asked (Ch.\,2, T).}
\end{qbox}
\begin{outline}
\begin{enumerate}
\item Circular: $r$, $v$ constant; acceleration in the rotating frame; gravity balances the centrifugal term.
\item Geosynchronous: $T=T_{sid}$, $a=42164$\,km (see \qref{A4}).
\item Ballistic: elliptic arcs with periapsis inside the Earth (ICBM).
\item Rectilinear: $h=0$, $e=1$, $p=0$; elliptic, parabolic, hyperbolic types.
\end{enumerate}
\end{outline}
\answer
\paragraph{Circular orbits.} $r=p=R$ and $v=\sqrt{\mu/R}$ are constant. In general
$\ddot{\underline r}=\ddot r\hat r+2\underline\omega\times(\dot r\hat r)+\dot{\underline\omega}\times\underline r+\underline\omega\times(\underline\omega\times\underline r)$
(radial, Coriolis, Euler, centripetal terms). On a circular orbit $\dot r=\ddot r=0$ and $\underline\omega=\underline h/R^2$ is constant, so
\begin{equation}
\ddot r\,\hat r=\underline 0=-\frac{\mu}{r^2}\hat r-\underline\omega\times(\underline\omega\times\underline r):
\end{equation}
the gravitational acceleration and the centrifugal (fictitious) acceleration balance at all times; the rotating frame attached to the
vehicle has no radial acceleration ($\omega^2R=\mu/R^2\Rightarrow v=\sqrt{\mu/R}$).

\paragraph{Geosynchronous orbits.} Period equal to the sidereal day (86164\,s): $a=42164$\,km, any $e$, $i$ (\qref{A4}).

\paragraph{Ballistic trajectories.} Arcs of ellipses whose periapsis lies \emph{inside} the Earth (apoapsis outside): the trajectory intersects
the surface twice. Used by ballistic missiles (ICBM) and sub-orbital flights.

\paragraph{Rectilinear trajectories.} $\underline v\parallel\underline r$ $\Rightarrow$ $\underline h=\underline 0$, $e=|-\hat r|=1$, $p=0$, $r_p=0$: motion along a line through the
attracting body ($\theta_*=\pm\pi$), non-uniform. Three types by energy: $a>0$, $\mathcal E<0$ \emph{rectilinear ellipse} (segment, falls back);
$a\to\infty$, $\mathcal E=0$ \emph{rectilinear parabola} ($v\to0$ at infinity); $a<0$, $\mathcal E>0$ \emph{rectilinear hyperbola} ($v$ does not vanish).
Example: vertical launch with no horizontal velocity.
\begin{keyf}
Circular: $\mu/R^2=\omega^2R$, $v=\sqrt{\mu/R}$; \ ballistic: $r_p<R_E$; \ rectilinear: $h=0$, $e=1$, type set by $\mathcal E=-\mu/2a$.
\end{keyf}
\pitfalls
\begin{itemize}
\item In the rotating frame the centrifugal term is fictitious; in the inertial frame gravity provides the centripetal acceleration.
\end{itemize}
\source{cap02.tex l.794--848 (circular), l.850--868 (geosynchronous), l.870--879 (ballistic), l.881--891 (rectilinear). Formulario: \fml{2-VCI}, \fml{2-BAL}.}

% ---------------------------------------------------------------------
\question{E6}{ECI and LVLH frames; orbit elements}
% source: cap02.tex l.1007-1156
\begin{qbox}
Define the Earth-centred inertial frame and the rotating LVLH frame, and the six orbit elements, including the special cases where some are undefined.
\qmeta{never asked (Ch.\,2, T\&N).}
\end{qbox}
\begin{outline}
\begin{enumerate}
\item ECI: $\hat c_3$ Earth axis (North), $\hat c_1$ vernal axis = ecliptic $\cap$ equator, $\hat c_2$ completes; precession of equinoxes neglected.
\item Orbit plane: $\Omega\in[-\pi,\pi[$, $i\in[0,\pi]$; $\omega$ from the node; $\theta_t=\omega+\theta_*$.
\item LVLH $(\hat r,\hat\theta,\hat h)=\underline{\underline R}_3(\theta_t)\underline{\underline R}_1(i)\underline{\underline R}_3(\Omega)(\hat c)$.
\item Elements $(a,e,i,\Omega,\omega)$ + one anomaly; special cases (circular, equatorial).
\end{enumerate}
\end{outline}
\answer
\paragraph{ECI frame.} Centred at the Earth's centre. $\hat c_3$ = Earth's rotation axis (towards the North pole); $\hat c_1$ = \textbf{vernal axis},
intersection of the ecliptic plane (plane of the Earth's orbit, normal to the ecliptic pole $\hat p$) with the equatorial plane; $\hat c_2=\hat c_3\times\hat c_1$.
Strictly, $\hat c_3$ precesses around $\hat p$ (precession of the equinoxes, $\sim$25\,700 years), but for satellite motion the axes are inertial.

\paragraph{Angles.} The orbit plane is identified by $\Omega$ (right ascension of the ascending node, $[-\pi,\pi[$, from $\hat c_1$ to the node $N$
where the spacecraft crosses the equator from South to North) and $i$ (inclination, $[0,\pi]$; direct $<90^\circ$, retrograde $>90^\circ$).
The \textbf{argument of perigee} $\omega\in[-\pi,\pi[$ is measured counter-clockwise from the nodal line to the periapsis; the argument of latitude
is $\theta_t=\omega+\theta_*$. Since $\underline h$ and $\underline e$ are first integrals, $(\Omega,i,\omega)$ are constant for Keplerian orbits.

\paragraph{LVLH frame} $(\hat r,\hat\theta,\hat h)$: radial, horizontal (in-plane, along the motion), normal to the orbit; it rotates with the spacecraft
(non-inertial) and is obtained from ECI with three counter-clockwise rotations:
\begin{equation}
\begin{bmatrix}\hat r\\ \hat\theta\\ \hat h\end{bmatrix}=\underline{\underline R}_3(\theta_t)\,\underline{\underline R}_1(i)\,\underline{\underline R}_3(\Omega)\begin{bmatrix}\hat c_1\\ \hat c_2\\ \hat c_3\end{bmatrix}.
\end{equation}

\paragraph{Orbit elements.} $a$ (size, energy), $e$ (shape), $(\Omega,i)$ (orientation of the plane), $\omega$ (periapsis in the plane): five constants
(five independent components of $\underline h$, $\underline e$); the sixth, time-varying, is one anomaly ($\theta_*$, $E$ or $M$), equivalent to the time of
periapsis passage. Special cases:
\begin{itemize}
\item circular ($e=0$): no periapsis, only $\theta_t$ is defined;
\item equatorial ($i=0,\pi$): no node, only the longitude of periapsis $\tilde\omega=\Omega\pm\omega$;
\item circular equatorial: only the true longitude $\tilde\lambda=\theta_t+\Omega$ ($i=0$) or $\Omega-\theta_t$ ($i=\pi$).
\end{itemize}
\begin{keyf}
$(\hat r,\hat\theta,\hat h)^T=\underline{\underline R}_3(\theta_t)\underline{\underline R}_1(i)\underline{\underline R}_3(\Omega)(\hat c_1,\hat c_2,\hat c_3)^T$, \ $\theta_t=\omega+\theta_*$; \
elements $(a,e,i,\Omega,\omega,\theta_*)$; $i\in[0,\pi]$, $\Omega,\omega\in[-\pi,\pi[$.
\end{keyf}
\pitfalls
\begin{itemize}
\item The ascending node is the South$\to$North crossing; $\omega$ is measured in the orbit plane from the node.
\end{itemize}
\source{cap02.tex l.1007--1099 (frames), l.1101--1156 (orbit elements). Formulario: \fml{2-OEL}, \fml{2-RA}.}

% ---------------------------------------------------------------------
\question{E7}{Ground track: latitude limits, symmetries, motion along the track}
% source: cap02.tex l.1657-1855
\begin{qbox}
Define the ground track and discuss its latitude limits, symmetry properties and the direction of motion along it.
\qmeta{never asked as a whole (Ch.\,2, T\&N); geosynchronous tracks in \qref{A4}.}
\end{qbox}
\begin{outline}
\begin{enumerate}
\item Definition; $\phi$ and $\lambda_g=\xi-\theta_g$ from the orbit elements.
\item Latitude limits: $\phi_{\max}=i$ (direct) or $\pi-i$ (retrograde); polar orbits.
\item Symmetries: $\omega=\pm\pi/2$ (meridian), $\omega=0,\pi$ (equator crossings), $e=0$ both.
\item Motion: retrograde always westward; direct depends on $a$, $e$, $\omega$; drift east if $T<T_{sid}$, west if $T>T_{sid}$.
\end{enumerate}
\end{outline}
\answer
\paragraph{Definition.} The ground track is the locus of the sub-satellite points on the rotating Earth, on a Mercator map ($\lambda_g$, $\phi$).
Equating the position written with orbit elements and with absolute longitude $\xi$ and latitude $\phi$:
\begin{equation}
s_\phi=s_{\theta_t}s_i,\qquad c_\xi=\frac{c_{\theta_t}c_\Omega-c_is_{\theta_t}s_\Omega}{c_\phi},\ \ s_\xi=\frac{c_{\theta_t}s_\Omega+c_is_{\theta_t}c_\Omega}{c_\phi},\qquad
\lambda_g=\xi-\theta_g,\ \ \theta_g=\theta_{g0}+\omega_E(t-t_0).
\end{equation}
\paragraph{Latitude limits.} $\phi=\arcsin(s_{\theta_t}s_i)$: extremes at $\theta_t=\pm\pi/2$. Direct orbits: $-i\le\phi\le i$; retrograde: $|\phi|\le\pi-i$
(latitude is in $[-\pi/2,\pi/2]$, inclination in $[0,\pi]$). Polar orbits ($i=\pi/2$) fly over all latitudes; crossing a pole the longitude jumps by
$\pi$ (Mercator discontinuity).
\paragraph{Symmetries} (from the symmetry of $\underline v$ w.r.t.\ the apsidal line and because points at the same latitude have the same eastward
inertial velocity): (a) $e\neq0$, $\omega=\pm\pi/2$: symmetric w.r.t.\ the meridian through the maximum/minimum latitude (apoapsis/periapsis);
(b) $e\neq0$, $\omega=0,\pi$: symmetric w.r.t.\ the equator-crossing points; (c) $e=0$: both. They allow reading $e$ and $\omega$ from a track.
\paragraph{Motion along the track.} \emph{Retrograde} orbits: the inertial velocity has a westward component, the track is always travelled East$\to$West.
\emph{Direct} orbits: the sub-satellite eastward velocity must be compared with that of the Earth at the same latitude: the track can be travelled
either way, depending on $a$, $e$, $\omega$. After one period the track point $P_1$ (same latitude and slope as $P_0$) moves east if
$T<T_{sid}$ (low orbits drift eastward) and west if $T>T_{sid}$ (high orbits move globally westward); geosynchronous direct orbits have closed
tracks, travelled partly westward and partly eastward.
\begin{keyf}
$s_\phi=s_{\theta_t}s_i$, $\lambda_g=\xi-\theta_g$; \ $\phi_{\max}=i$ (direct), $\pi-i$ (retrograde); \ drift east if $T<T_{sid}$, west if $T>T_{sid}$.
\end{keyf}
\pitfalls
\begin{itemize}
\item Retrograde: $\phi_{\max}=\pi-i$, not $i$.
\end{itemize}
\source{cap02.tex l.1657--1739 (ground track), l.1741--1796 (latitude, symmetry), l.1845--1855 (motion along the track). Formulario: \fml{2-LAT}, \fml{2-GTS}, \fml{2-GTD}.}

% ---------------------------------------------------------------------
\question{E8}{Three-dimensional orbit transfers}
% source: cap03.tex l.451-565
\begin{qbox}
Discuss in-plane and out-of-plane velocity changes in three-dimensional orbit transfers and describe a near-optimal LEO--GEO transfer.
\qmeta{never asked (Ch.\,3, T; LEO--GEO is N).}
\end{qbox}
\begin{outline}
\begin{enumerate}
\item In-plane: $dv=\frac{\mu}{2a^2v}da$ $\Rightarrow$ most efficient at periapsis (max $v$).
\item Out-of-plane: rotation by $\Delta\phi$ about $BP$, $|\Delta\underline v|=2v\,s_{\Delta\phi/2}$ $\Rightarrow$ at apoapsis (min $v$), at a node.
\item LEO--GEO: (A) 3 impulses; (B) 2 impulses combining circularisation and plane change; true optimum splits the plane change.
\item Reversibility of optimal transfers.
\end{enumerate}
\end{outline}
\answer
\paragraph{In-plane velocity changes.} From $\mathcal E=-\frac{\mu}{2a}=\frac{v^2}{2}-\frac{\mu}{r}$ at fixed $r$: $\frac{\mu}{2a^2}da=v\,dv$, i.e.\
$dv=\frac{\mu}{2a^2v}da$: the $\Delta v$ needed for a given change of $a$ is minimum where $v$ is maximum, at \textbf{periapsis}.
\paragraph{Out-of-plane velocity changes.} To rotate the orbit plane by $\Delta\phi$ about the line $BP$ ($B$ = attracting body) keeping the speed:
$|\underline v_0|=|\underline v_f|=v$ and
\begin{equation}
|\Delta\underline v|=2v\sin\frac{\Delta\phi}{2},
\end{equation}
expensive (e.g.\ $\Delta\phi=60^\circ$ costs $\Delta v=v$), so it is convenient where $v$ is minimum, at \textbf{apoapsis}; to change $i$ without changing
$\Omega$ the impulse must be at a node.
\paragraph{Near-optimal LEO--GEO transfer.} From a circular LEO ($R_i$, $i_0>0$) to a circular GEO ($R_f$, $i_f<i_0$):
(A) \emph{3 impulses}: at perigee raise $r_A$ to $R_f$; at apoapsis circularise; at a node rotate the plane by $i_0-i_f$.
(B) \emph{2 impulses}: first impulse at a node to raise $r_A=R_f$, second at the opposite node (the apoapsis) combining circularisation and
plane change: $\Delta v_2^{(B)}=\sqrt{v_A^2+v_{cf}^2-2v_Av_{cf}c_{\Delta\phi}}<\Delta v_2^{(A)}+\Delta v_3^{(A)}$, with
$v_A=\sqrt{\frac{2\mu}{R_i+R_f}\frac{R_i}{R_f}}$, $v_{cf}=\sqrt{\mu/R_f}$ (vector addition is cheaper than two separate impulses).
(B) is near-optimal; the true optimum also performs a small part of the plane change at the first impulse
($\Delta v_1=\sqrt{v_{ci}^2+v_{PH}^2-2v_{ci}v_{PH}c_{(i_1-i_H)}}$).
\paragraph{Remark.} If the transfer A$\to$B is optimal, the return transfer B$\to$A, with impulses of equal magnitude and opposite direction, is optimal
too (e.g.\ interior Hohmann transfer).
\begin{keyf}
$dv=\dfrac{\mu}{2a^2v}da$; \ $|\Delta v|=2v\sin\dfrac{\Delta\phi}{2}$; \ combined: $\Delta v=\sqrt{v_1^2+v_2^2-2v_1v_2\cos\Delta\phi}$.
\end{keyf}
\pitfalls
\begin{itemize}
\item In-plane at periapsis, plane changes at apoapsis (and at a node).
\end{itemize}
\source{cap03.tex l.451--537 (3D transfers, LEO--GEO), l.539--565 (concluding remarks). Formulario: \fml{3-3DV}, \fml{3-LGE}.}

% ---------------------------------------------------------------------
\question{E9}{Sphere of influence and qualitative effect of a flyby}
% source: cap04.tex l.20-126, l.522-600
\begin{qbox}
Introduce the patched-conic approximation and derive the radius of the sphere of influence; describe qualitatively the effect of a planetary flyby.
\qmeta{never asked (Ch.\,4: introduction T, SoI T\&N, qualitative flyby T).}
\end{qbox}
\begin{outline}
\begin{enumerate}
\item Patched conics: one attracting body at a time; departure hyperbola, heliocentric ellipse, arrival hyperbola.
\item Equations of the vehicle relative to the planet and to the Sun: dominant + perturbing terms.
\item SoI: equal ratios $\Rightarrow$ $\rho_{SoI}=L(m_p/m_s)^{2/5}$; values.
\item Flyby: $\underline v^{(H)}_\pm=\underline v^{(H)}_p+\underline v_\infty^\pm$; $|\underline v_\infty|$ unchanged, direction rotated $\Rightarrow$ heliocentric speed and energy change.
\item Behind the planet: increase; in front: decrease; clockwise/counter-clockwise.
\end{enumerate}
\end{outline}
\answer
\paragraph{Patched conics.} The spacecraft is assumed to be subject to \textbf{one attracting body at a time}: near the departure planet only its gravity
(hyperbola), far away only the Sun (heliocentric ellipse), near the arrival planet only its gravity (hyperbola). Each arc is Keplerian.

\paragraph{Sphere of influence.} For Sun ($S$), planet ($P$), vehicle ($V$), with $\underline\rho$ = vehicle w.r.t.\ planet, $\underline r$ = vehicle w.r.t.\ Sun,
$\underline L$ = planet w.r.t.\ Sun:
\begin{equation}
\ddot{\underline\rho}=-\frac{Gm_p}{\rho^3}\underline\rho-Gm_s\left[\frac{\underline r}{r^3}-\frac{\underline L}{L^3}\right],\qquad
\ddot{\underline r}=-\frac{Gm_s}{r^3}\underline r-Gm_p\frac{\underline\rho}{\rho^3}.
\end{equation}
In each equation the first term is the dominant attraction, the second the perturbation. The SoI is where the perturbation/dominant ratios are
equal:
$\frac{m_s|\underline r/r^3-\underline L/L^3|}{m_p/\rho^2}=\frac{m_p/\rho^2}{m_s/r^2}$. With $\frac{\underline r}{r^3}-\frac{\underline L}{L^3}\simeq\frac{\underline\rho}{L^3}$ and $r\simeq L$:
\begin{equation}
\left(\frac{\rho}{L}\right)^{5}=\left(\frac{m_p}{m_s}\right)^2\;\Longrightarrow\;\rho_{SoI}=L\left(\frac{m_p}{m_s}\right)^{2/5}.
\end{equation}
An estimate, not an exact boundary: Earth $9.24\cdot10^5$\,km, Mars $5.74\cdot10^5$\,km, Jupiter $4.83\cdot10^7$\,km, Neptune $8.67\cdot10^7$\,km (larger
than Jupiter's because it is farther from the Sun).

\paragraph{Flyby (qualitative).} Inside the SoI, with no manoeuvre, the spacecraft follows a hyperbola: the relative velocity at the SoI keeps its
magnitude $v_\infty$ but is rotated by the deflection angle. In the heliocentric frame
\begin{equation}
\underline v^{(H)}_\pm=\underline v_p^{(H)}+\underline v_\infty^\pm ,
\end{equation}
so the heliocentric velocity changes in magnitude and direction and the heliocentric energy changes ($\mathcal E_+\neq\mathcal E_-$). The energy is
exchanged with the planet, whose velocity change is negligible. Depending on the geometry a flyby can increase, leave unchanged or reduce the
heliocentric speed: \textbf{periapsis behind the planet} (passing behind it w.r.t.\ its motion) $\Rightarrow$ \textbf{increase}; hyperbola symmetric
w.r.t.\ $\underline v_p^{(H)}$ $\Rightarrow$ unchanged; \textbf{periapsis in front} $\Rightarrow$ \textbf{decrease}; this holds for both clockwise and counter-clockwise
flybys. Gravity assists made Voyager, Cassini, Ulysses (out of the ecliptic) feasible.
\begin{keyf}
$\rho_{SoI}=L\left(\dfrac{m_p}{m_s}\right)^{2/5}$; \ $\underline v^{(H)}_\pm=\underline v_p^{(H)}+\underline v_\infty^\pm$, $|\underline v_\infty^+|=|\underline v_\infty^-|$; \ behind $\Rightarrow$ speed up, in front $\Rightarrow$ slow down.
\end{keyf}
\pitfalls
\begin{itemize}
\item The speed relative to the planet does not change; the heliocentric one does.
\end{itemize}
\source{cap04.tex l.20--27 (patched conics), l.29--126 (sphere of influence), l.522--600 (flyby). Formulario: \fml{4-SOI}, \fml{4-FBC}.}

% ---------------------------------------------------------------------
\question{E10}{Problem of two bodies: relative and absolute motion}
% source: cap05.tex l.111-242
\begin{qbox}
Describe the motion of two bodies subject to their mutual attraction: relative motion, motion of the centre of mass and absolute motions.
\qmeta{never asked (Ch.\,5, T).}
\end{qbox}
\begin{outline}
\begin{enumerate}
\item Relative motion: Keplerian with $G(m_1+m_2)$.
\item Centre of mass on the line of the bodies, uniform motion; origin $O'$ there.
\item Absolute motions about $O'$: Keplerian with equivalent masses $m_2^3/M^2$, $m_1^3/M^2$.
\item Same eccentricity, $a_1/a_2=m_2/m_1$, $\omega^2=GM/(a_1+a_2)^3$; planet--satellite limit.
\end{enumerate}
\end{outline}
\answer
\paragraph{Relative motion.} From $m_1\ddot{\underline R}_1=\underline F_1$, $m_2\ddot{\underline R}_2=\underline F_2=-\underline F_1$:
$\frac{d^2}{dt^2}(\underline R_2-\underline R_1)=-\frac{G(m_1+m_2)}{r_{12}^3}(\underline R_2-\underline R_1)$: 2 moves relative to 1 as if the total mass
$M=m_1+m_2$ were at 1: Keplerian.
\paragraph{Centre of mass.} $\underline R_c=\frac{m_1\underline R_1+m_2\underline R_2}{M}$ lies on the segment joining the bodies
($\underline R_1-\underline R_c=\frac{m_2}{M}(\underline R_1-\underline R_2)$, opposite to $\underline R_2-\underline R_c$) and moves uniformly; take it as origin $O'$:
$m_1\underline R_1+m_2\underline R_2=\underline 0$, $m_1R_1=m_2R_2$.
\paragraph{Absolute motion.} With $r_{12}=R_1(1+m_1/m_2)$:
\begin{equation}
\ddot{\underline R}_1=-G\frac{m_2^3}{M^2}\frac{\underline R_1}{R_1^3},\qquad \ddot{\underline R}_2=-G\frac{m_1^3}{M^2}\frac{\underline R_2}{R_2^3}:
\end{equation}
each body moves on a Keplerian conic about $O'$ with equivalent mass $m_2^3/M^2$ (resp.\ $m_1^3/M^2$). Since $R_1/R_2=m_2/m_1$ is constant, the bodies
reach periapsis and apoapsis together: $\frac{a_1(1-e_1)}{a_2(1-e_2)}=\frac{a_1(1+e_1)}{a_2(1+e_2)}\Rightarrow e_1=e_2$, and $a_1/a_2=m_2/m_1$. The larger
mass moves entirely inside the orbit of the smaller one if $e<(m_1-m_2)/M$. Mean motion: $\omega^2=\frac{GM}{(a_1+a_2)^3}$.
\paragraph{Planet--satellite problem.} $m_1\gg m_2$: the centre of mass coincides with the planet's centre and the satellite obeys the classical
one-body equation (restricted two-body problem).
\begin{keyf}
Relative: $\mu=G(m_1+m_2)$; \ $m_1\underline R_1+m_2\underline R_2=\underline 0$; \ absolute: equivalent masses $m_2^3/M^2$, $m_1^3/M^2$; \ $e_1=e_2$, $\frac{a_1}{a_2}=\frac{m_2}{m_1}$, $\omega^2=\frac{GM}{(a_1+a_2)^3}$.
\end{keyf}
\pitfalls
\begin{itemize}
\item Relative and absolute motions are both Keplerian, but with different gravitational parameters.
\end{itemize}
\source{cap05.tex l.111--236 (problem of two bodies), l.238--242 (planet--satellite).}

% ---------------------------------------------------------------------
\question{E11}{Newton's law with propulsion: thrust and effective exhaust velocity}
% source: cap06.tex l.21-110
\begin{qbox}
Derive the equation of motion of a rocket (Newton's law with propulsion) and define thrust, effective exhaust velocity and specific impulse.
\qmeta{never asked (Ch.\,6, T).}
\end{qbox}
\begin{outline}
\begin{enumerate}
\item Momentum of vehicle + ejected mass at $t$ and $t+\Delta t$; external forces $\underline G+\underline A+\underline P_{ext}$, $\underline P_{ext}=(p_e-p_a)A_e\hat v$.
\item Limit $\Delta t\to0$: $m\dot{\underline v}+\dot{\tilde m}\underline c_{ej}=\underline G+\underline A+A_e(p_e-p_a)\hat v$.
\item $\underline T=\dot m\underline c$, $\underline c=\underline c_{ej}+A_e(p_e-p_a)\hat v/\dot m$, $\dot m=-T/c$; $\dot{\underline v}=(\underline G+\underline A+\underline T)/m$.
\item $c=g_0I_{sp}$.
\end{enumerate}
\end{outline}
\answer
\paragraph{Momentum balance.} At $t$ the vehicle has mass $m$ and velocity $\underline v$; after $\Delta t$ the ejected mass $\Delta\tilde m$ moves with
$\underline v+\underline c_{ej}$ ($\underline c_{ej}$ = exhaust velocity relative to the vehicle) and the vehicle has mass $m-\Delta\tilde m$ and velocity $\underline v+\Delta\underline v$:
\begin{equation}
\Delta\underline M=(m-\Delta\tilde m)(\underline v+\Delta\underline v)+\Delta\tilde m(\underline v+\underline c_{ej})-m\underline v .
\end{equation}
External forces: gravity $\underline G$, aerodynamic force $\underline A$ and the external pressure force; on the nozzle exit (area $A_e$, pressure $p_e$) vs
ambient pressure $p_a$, the net pressure force is $\underline P_{ext}=(p_e-p_a)A_e\hat v$. Dividing by $\Delta t\to0$ (neglecting $\Delta\tilde m\Delta\underline v$):
\begin{equation}
m\frac{d\underline v}{dt}+\dot{\tilde m}\,\underline c_{ej}=\underline G+\underline A+A_e(p_e-p_a)\hat v,\qquad \dot{\tilde m}=-\dot m .
\end{equation}
\paragraph{Thrust.} Hence $\frac{d\underline v}{dt}=\frac{\underline G+\underline A+\underline T}{m}$ with
\begin{equation}
\underline T=\dot m\,\underline c_{ej}+A_e(p_e-p_a)\hat v=\dot m\,\underline c,\qquad \underline c=\underline c_{ej}+\frac{A_e(p_e-p_a)}{\dot m}\hat v,\qquad \dot m=-\frac{T}{c}<0,
\end{equation}
$\underline c$ = \textbf{effective exhaust velocity} (includes the pressure term); $\underline T$ is directed opposite to the ejected mass. $c$ is expressed through the
\textbf{specific impulse} $I_{sp}$ [s]: $c=g_0I_{sp}$ ($g_0$ = sea-level gravity). These equations, with $\dot{\underline r}=\underline v$, are the equations of spaceflight
whose controls are the thrust magnitude and direction.
\begin{keyf}
$\dot{\underline v}=\dfrac{\underline G+\underline A+\underline T}{m}$, \ $\underline T=\dot m\underline c$, \ $\underline c=\underline c_{ej}+\dfrac{A_e(p_e-p_a)}{\dot m}\hat v$, \ $\dot m=-\dfrac Tc$, \ $c=g_0I_{sp}$.
\end{keyf}
\pitfalls
\begin{itemize}
\item $\dot m<0$: keep track of the sign between the mass lost and the vehicle mass.
\end{itemize}
\source{cap06.tex l.21--110 (Newton's law with propulsion), l.573--669 (equations of spaceflight). Formulario: \fml{6-ISP}.}

% ---------------------------------------------------------------------
\question{E12}{Optimal rocket staging}
% source: cap06.tex l.225-354
\begin{qbox}
Formulate the optimal staging problem and solve it for similar stages; comment on the maximum $\Delta V$ as a function of the number of stages.
\qmeta{never asked (Ch.\,6, T; the special-case solution is N).}
\end{qbox}
\begin{outline}
\begin{enumerate}
\item Maximise $\Delta v=-\sum c_i\ln u_0^{(i)}$ with the payload-ratio constraint $\frac{m_u^{(N)}}{m_0^{(1)}}=\prod\frac{u_0^{(i)}-\epsilon_s^{(i)}}{1-\epsilon_s^{(i)}}$.
\item Similar stages: Lagrange multiplier $\Rightarrow$ all $u_0^{(i)}$ equal.
\item $u_0=\epsilon_s+(1-\epsilon_s)[m_u^{(N)}/m_0^{(1)}]^{1/N}$, $\Delta v_{\max}=-cN\ln u_0$.
\item Finite limit for $N\to\infty$; diminishing returns; 3--4 stages.
\end{enumerate}
\end{outline}
\answer
\paragraph{Problem.} Given the structural coefficients $\epsilon_s^{(i)}$, the exhaust velocities $c_i$ and the payload ratio
$m_u^{(N)}/m_0^{(1)}$, find the mass ratios $u_0^{(i)}$ that maximise $\Delta v=-\sum_{i}c_i\ln u_0^{(i)}$, subject to
\[\frac{m_u^{(N)}}{m_0^{(1)}}=\prod_i\frac{m_u^{(i)}}{m_0^{(i)}}=\prod_{i=1}^N\frac{u_0^{(i)}-\epsilon_s^{(i)}}{1-\epsilon_s^{(i)}}.\]
\paragraph{Similar stages} ($\epsilon_s^{(i)}=\epsilon_s$, $c_i=c$). Minimise
$H=c\sum\ln u_0^{(i)}+\lambda\left[\frac{m_u^{(N)}}{m_0^{(1)}}-\frac{\prod(u_0^{(i)}-\epsilon_s)}{(1-\epsilon_s)^N}\right]$: $\partial H/\partial\lambda=0$ gives the constraint,
$\partial H/\partial u_0^{(i)}=0$ gives $\frac{c}{u_0^{(i)}}=\frac{\lambda}{(1-\epsilon_s)^N}\prod_{j\neq i}(u_0^{(j)}-\epsilon_s)$ for every $i$; dividing two of them,
$\frac{u_0^{(k)}-\epsilon_s}{u_0^{(k)}}=\frac{u_0^{(1)}-\epsilon_s}{u_0^{(1)}}$ $\Rightarrow$ \textbf{all stages have the same $u_0$}. From the constraint:
\begin{equation}
u_0=\epsilon_s+(1-\epsilon_s)\left[\frac{m_u^{(N)}}{m_0^{(1)}}\right]^{1/N},\qquad
\Delta v_{\max}=-c\,N\ln\left\{\epsilon_s+(1-\epsilon_s)\left[\frac{m_u^{(N)}}{m_0^{(1)}}\right]^{1/N}\right\}.
\end{equation}
\paragraph{Comments.} $\Delta v_{\max}$ increases with $N$ but with diminishing returns and tends to a finite limit (an upper bound that can never be
exceeded), while structural and operational complexity grow: launchers rarely have more than 5 stages (often 3--4).
\begin{keyf}
equal $u_0$ for similar stages; \ $u_0=\epsilon_s+(1-\epsilon_s)(m_u^{(N)}/m_0^{(1)})^{1/N}$; \ $\Delta v_{\max}=-cN\ln u_0$.
\end{keyf}
\pitfalls
\begin{itemize}
\item The exponent $1/N$ applies only to the payload ratio (source errata \#6).
\end{itemize}
\source{cap06.tex l.225--354 (optimal staging; errata \#5--6). Formulario: \fml{6-OPT}.}

% ---------------------------------------------------------------------
\question{E13}{Ascent trajectory of a launch vehicle}
% source: cap06.tex l.356-502
\begin{qbox}
Describe the phases of the ascent trajectory of a launch vehicle, the selection of the orbit plane and the evolution of the osculating elements.
\qmeta{never asked (Ch.\,6, T).}
\end{qbox}
\begin{outline}
\begin{enumerate}
\item Atmospheric arc ($<40$\,km) and exoatmospheric arc.
\item Phases: vertical arc, roll, pitch, gravity turn, rarefied-atmosphere arc, coast, injection.
\item Launch-site equation $\cos i=\cos\phi_L\sin\psi_L$; $i_{\min}=|\phi_L|$ (East); $\Omega$ from launch time.
\item Osculating elements: $\Omega$, $i$ early; $a$, $e$ at injection; $e\simeq1$ at lift-off.
\end{enumerate}
\end{outline}
\answer
\paragraph{Two arcs.} Atmospheric arc (non-negligible density, below $\sim40$\,km: structural integrity must be preserved) and exoatmospheric arc.
\paragraph{Phases.}
(a) \textbf{vertical arc}: launch-site safety (exhaust directed vertically) and to avoid a premature rotation of $\underline v$; (b) \textbf{roll manoeuvre}: to assume
the attitude for the pitch; (c) \textbf{pitch manoeuvre}: rotates $\underline v$ to the desired flight-path and heading angles, completing the orbit-plane
selection; (d) \textbf{gravity turn}: thrust (longitudinal axis) aligned with the velocity relative to the atmosphere ($\alpha=0$) to avoid destructive bending
loads, gravity turns the trajectory (first/second stage); some vehicles tolerate a small incidence with $q\alpha\le(q\alpha)_{\max}$, $q=\frac12\rho v_R^2$;
(e) \textbf{rarefied-atmosphere arc}: incidence allowed within the $q\alpha$ limit, ends at burnout of the 2nd/3rd stage; (f) \textbf{coast arc}: ballistic,
nearly Keplerian, altitude increases, speed decreases; (g) \textbf{injection arc}: upper stage burn near the target altitude raises the speed to the orbital
value at nearly constant altitude. Initial velocity $v_0=R_E\omega_E\cos\phi_L$ (Earth's rotation), final $v_f=\sqrt{\mu_E/R_f}$.
\paragraph{Orbit plane.} Spherical trigonometry gives the \textbf{equation of the launch site}: $\cos i=\cos\phi_L\sin\psi_L$ ($\psi_L$ = launch azimuth from North).
Minimum inclination $i_{\min}=|\phi_L|$ launching East ($\psi_L=\pi/2$), maximum $\pi-|\phi_L|$ launching West (rare: the Earth's rotation helps eastward).
Hence $|\phi_L|\le i\le\pi-|\phi_L|$. $\Omega$ is selected with the launch date and time (at most two opportunities per day).
\paragraph{Osculating elements} (those of the Keplerian orbit the vehicle would follow if thrust and perturbations ceased): $\Omega$ and $i$ change in the
first phases and then remain almost constant (plane selected early, by roll and pitch); $a$ and $e$ change mostly in the last phases, especially
during injection, which circularises the orbit. Right after launch $e\simeq1$ (nearly rectilinear path).
\begin{keyf}
$\cos i=\cos\phi_L\sin\psi_L$; \ $|\phi_L|\le i\le\pi-|\phi_L|$; \ $v_0=R_E\omega_E\cos\phi_L$; \ gravity turn: $\alpha=0$; \ $q\alpha\le(q\alpha)_{\max}$.
\end{keyf}
\pitfalls
\begin{itemize}
\item A site cannot reach inclinations below its latitude without a plane change.
\end{itemize}
\source{cap06.tex l.356--482 (phases, launch-site equation), l.484--502 (osculating elements). See also \qref{A9} (losses).}

% ---------------------------------------------------------------------
\question{E14}{Orbit perturbations: overview and Lagrange planetary equations}
% source: cap08.tex l.41-118, l.338-372
\begin{qbox}
List the main orbit perturbations of an Earth satellite, their nature and relative importance at different altitudes, and derive the vector equations
for $\underline h$ and $\underline e$ in perturbed motion.
\qmeta{never asked as such; possibly the generic "Perturbations" of Mar 2025 (Ch.\,8, T; equations "outline").}
\end{qbox}
\begin{outline}
\begin{enumerate}
\item Perturbations: asphericity, third body (conservative); drag, SRP (surface); altitude dependence.
\item Importance vs altitude: LEO, MEO, HEO, GEO; magnitudes at 1000\,km.
\item $\ddot{\underline r}=-\frac{\mu}{r^3}\underline r+\underline f$: $\dot{\underline h}=\underline r\times\underline f$, $\dot{\underline e}=\frac1\mu[\underline f\times\underline h+\underline v\times(\underline r\times\underline f)]$.
\item Osculating elements; Gauss equations (components $f_r,f_\theta,f_h$); which components change what.
\end{enumerate}
\end{outline}
\answer
\paragraph{Perturbations.} (a) Earth asphericity (mainly $J_2$, oblateness), (b) third-body attraction (Sun, Moon), (c) aerodynamic drag, (d) solar radiation
pressure. (a) and (b) are conservative; (c) and (d) are surface perturbations (depend on area/mass). All except (d) depend on the altitude.
\paragraph{Importance vs altitude.} LEO (up to 700\,km): oblateness and drag dominate. MEO (700 to 10\,000\,km): oblateness dominates; SRP and
third body are non-negligible; drag is negligible above about 1000\,km. High orbits ($>$10\,000\,km): SRP, third body and oblateness (decreasing fast). GEO:
$J_{22}$ dominates by resonance. At 1000\,km ($a_0=\mu/r^2$):
$a_{J2}\simeq10^{-2}a_0$, $a_{3B}\simeq10^{-7}a_0$, $a_{RP}\simeq10^{-9}a_0$, $a_D\simeq10^{-10}a_0$.
\paragraph{Vector equations.} With $\ddot{\underline r}=-\frac{\mu}{r^3}\underline r+\underline f$ ($\underline f$ may include thrust), $\underline h$ and $\underline e$ are no longer constant:
\begin{equation}
\frac{d\underline h}{dt}=\underline v\times\underline v+\underline r\times\left[-\frac{\mu}{r^3}\underline r+\underline f\right]=\underline r\times\underline f,\qquad
\frac{d\underline e}{dt}=\frac1\mu\left[\underline f\times\underline h+\underline v\times(\underline r\times\underline f)\right],
\end{equation}
the second because the Keplerian terms cancel exactly as in the unperturbed proof. They are frame-independent; projected on $(\hat r,\hat\theta,\hat h)$
they give the Gauss (Lagrange planetary) equations for the \textbf{osculating elements} (those of the Keplerian orbit tangent to the actual one at time $t$), e.g.\
$\dot p=2\sqrt{p/\mu}\,rf_\theta$, $\dot i=\frac{rf_hc_{\theta_t}}{h}$, $\dot\Omega=\frac{rf_hs_{\theta_t}}{hs_i}$,
$\dot a=\frac{2a^2}{h}(es_{\theta_*}f_r+\frac prf_\theta)$. In-plane components ($f_r,f_\theta$) change $a$, $e$, $\omega$; only $f_h$ changes the plane
($i$, $\Omega$). Singularities for $e\to0$ ($\omega$, $\theta_*$) and $i\to0$ ($\Omega$, $\omega$) disappear using $\theta_t$ or $\Omega+\omega$.
\begin{keyf}
$\dot{\underline h}=\underline r\times\underline f$, \ $\dot{\underline e}=\frac1\mu[\underline f\times\underline h+\underline v\times(\underline r\times\underline f)]$; \ $f_h$ $\Rightarrow$ $i,\Omega$; \ $f_r,f_\theta$ $\Rightarrow$ $a,e,\omega$.
\end{keyf}
\pitfalls
\begin{itemize}
\item The osculating elements change; they describe the instantaneous Keplerian orbit.
\end{itemize}
\source{cap08.tex l.41--85 (introduction, perturbations on different orbits), l.87--116 (vector equations), l.338--372 (Gauss equations; errata \#10). Formulario: \fml{8-EOM}, \fml{8-GAU}.}

% ---------------------------------------------------------------------
\question{E15}{Euler's principal rotation theorem}
% source: cap09.tex l.426-680
\begin{qbox}
State Euler's principal rotation theorem and the main remarks on principal axis and angle.
\qmeta{never asked (Ch.\,9: only the statement may be requested; remarks T\&N).}
\end{qbox}
\begin{outline}
\begin{enumerate}
\item Statement; eigen-axis: eigenvector of $\underline{\underline R}$ with eigenvalue 1, same components in both frames.
\item $\underline{\underline R}=c_\Phi\underline{\underline I}+(1-c_\Phi)\underline a\underline a^T-s_\Phi\tilde a$; $\Phi=\arccos\frac{r_{11}+r_{22}+r_{33}-1}{2}$, $a_i$ from the skew part.
\item Remarks: $(\underline a,\Phi)\equiv(-\underline a,-\Phi)$; $\Phi=2k\pi$ no axis; 4 parameters, 1 constraint.
\end{enumerate}
\end{outline}
\answer
\paragraph{Statement.} A rigid body (or frame) can be brought from an arbitrary initial orientation to an arbitrary final orientation by a
\textbf{single rigid rotation through a principal angle $\Phi$ about a principal axis $\hat a$}. The axis is fixed in both orientations (eigen-axis):
its components along $(\hat E_i)$ and $(\hat e_i)$ coincide, $(\underline{\underline R}^T-\underline{\underline I})\underline a=\underline 0$, i.e.\ $\underline a$ is the eigenvector of
$\underset{B\leftarrow N}{\underline{\underline R}}$ associated with the eigenvalue 1 (every proper orthogonal matrix has it).
\paragraph{Rotation matrix and inverse relations.}
\begin{equation}
\begin{gathered}
\underset{B\leftarrow N}{\underline{\underline R}}=c_\Phi\,\underline{\underline I}+(1-c_\Phi)\,\underline a\,\underline a^T-s_\Phi\,\tilde a,\\[2pt]
\Phi=\arccos\frac{r_{11}+r_{22}+r_{33}-1}{2}\in[0,\pi],\quad a_1=\frac{r_{23}-r_{32}}{2s_\Phi},\ a_2=\frac{r_{31}-r_{13}}{2s_\Phi},\ a_3=\frac{r_{12}-r_{21}}{2s_\Phi}
\end{gathered}
\end{equation}
($\Phi=\pi$: $a_j=\sqrt{(r_{jj}+1)/2}$ for the largest $r_{jj}$, then the off-diagonal terms).
\paragraph{Remarks.} (a) $(\underline a,\Phi)$ and $(-\underline a,-\Phi)$ give the same attitude; (b) also $(-\underline a,2k\pi-\Phi)$; (c) the axis is unique
($\pm\underline a$) except for $\Phi=2k\pi$ ($\underline{\underline R}=\underline{\underline I}$, no axis); (d) opposite rotations: $(-\underline a,\Phi)$ or $(\underline a,-\Phi)$;
(e) 4 parameters with 1 constraint ($|\underline a|=1$): redundancy 1. The theorem is the basis of the Euler parameters (\qref{D4}) and of eigen-axis manoeuvres.
\begin{keyf}
single rotation $\Phi$ about $\hat a$; $\underline{\underline R}\underline a=\underline a$; \ $\underline{\underline R}=c_\Phi\underline{\underline I}+(1-c_\Phi)\underline a\underline a^T-s_\Phi\tilde a$; \ $1+2c_\Phi=\mathrm{tr}\,\underline{\underline R}$.
\end{keyf}
\pitfalls
\begin{itemize}
\item The theorem is about the existence of a single rotation; in the theory part only the statement is required.
\end{itemize}
\source{cap09.tex l.426--466 (statement), l.468--619 (proofs, not required), l.621--680 ($\Phi$, $\underline a$ from $\underline{\underline R}$; remarks). Formulario: \fml{9-PAX}, \fml{9-PHI}, \fml{9-PRM}.}

% ---------------------------------------------------------------------
\question{E16}{Euler's equations of attitude dynamics and kinetic energy}
% source: cap11.tex l.21-58, l.386-473
\begin{qbox}
Derive Euler's equations of attitude dynamics and the expression of the rotational kinetic energy of a rigid body.
\qmeta{never asked (Ch.\,11, T\&N).}
\end{qbox}
\begin{outline}
\begin{enumerate}
\item $\underline H_O=\underline H_{MC}+\underline H_C$; $\dot{\underline H}_C=\underline L_C$ ($C$ centre of mass).
\item Transport theorem in body axes: $\underline{\underline I}\dot{\underline\omega}+\tilde\omega\underline{\underline I}\underline\omega=\underline L$; principal axes form.
\item $T_{rot}=\frac12\underline\omega\cdot\underline H_C=\frac12\underline\omega^T\underline{\underline I}\underline\omega$; $\dot T_{rot}=\underline\omega\cdot\underline L_C$.
\end{enumerate}
\end{outline}
\answer
\paragraph{Angular momentum.} With $\underline R=\underline R_C+\underline r$: $\underline H_O=\int\underline R\times\dot{\underline R}\,dm=M\underline R_C\times\dot{\underline R}_C+\int\underline r\times\dot{\underline r}\,dm$
(mixed terms vanish by definition of centre of mass): orbital part plus $\underline H_C$, the angular momentum about $C$. For a rigid body
$\dot{\underline r}=\underline\omega\times\underline r$ and $\underline H_C=\underrightarrow I_C\cdot\underline\omega$. If $C$ is the centre of mass (or a point with no inertial
velocity): $\dot{\underline H}_C=\underline L_C$.
\paragraph{Euler's equations.} Using the transport theorem with the body frame $B$, in which $\underline{\underline I}_C^{(B)}$ is constant:
\begin{equation}
\frac{{}^Bd\underline H_C}{dt}+\underline\omega\times\underline H_C=\underline L_C
\;\Longrightarrow\;
\underline{\underline I}_C^{(B)}\dot{\underline\omega}^{(B)}+\tilde\omega\,\underline{\underline I}_C^{(B)}\underline\omega^{(B)}=\underline L_C^{(B)}.
\end{equation}
In principal axes:
\begin{equation}
I_1\dot\omega_1=(I_2-I_3)\omega_2\omega_3+L_1,\quad I_2\dot\omega_2=(I_3-I_1)\omega_1\omega_3+L_2,\quad I_3\dot\omega_3=(I_1-I_2)\omega_1\omega_2+L_3 .
\end{equation}
They relate the external torque to the angular velocity (nonlinear, coupled); attitude follows from the kinematic equations.
\paragraph{Kinetic energy.} $T=\frac12M\dot R_C^2+\frac12\int\dot{\underline r}\cdot\dot{\underline r}\,dm=T_{transl}+T_{rot}$, and with $\dot{\underline r}=\underline\omega\times\underline r$:
\begin{equation}
T_{rot}=\frac12\int(\underline\omega\times\underline r)\cdot(\underline\omega\times\underline r)\,dm=\frac12\underline\omega\cdot\underline H_C=\frac12\underline\omega^{(B)T}\underline{\underline I}_C^{(B)}\underline\omega^{(B)},\qquad
\dot T_{rot}=\underline\omega\cdot\underline L_C .
\end{equation}
Work: $W=\int(\underline F\cdot\dot{\underline R}_C+\underline L_C\cdot\underline\omega)\,dt$. Without torque, $T_{rot}$ and $\underline H_C$ are constant.
\begin{keyf}
$\dot{\underline H}_C=\underline L_C$; \ $\underline{\underline I}\dot{\underline\omega}+\tilde\omega\underline{\underline I}\underline\omega=\underline L$; \ $I_1\dot\omega_1=(I_2-I_3)\omega_2\omega_3+L_1$ (cyclic); \ $T_{rot}=\frac12\underline\omega^T\underline{\underline I}\underline\omega$, $\dot T_{rot}=\underline\omega\cdot\underline L$.
\end{keyf}
\pitfalls
\begin{itemize}
\item Euler's equations hold about the centre of mass (or a fixed point), with components in body axes.
\end{itemize}
\source{cap11.tex l.21--58 (angular momentum), l.386--446 (Euler's equations), l.448--473 (kinetic energy). Formulario: \fml{11-EUL}, \fml{11-KEN}.}

% ---------------------------------------------------------------------
\question{E17}{Properties of the inertia matrix; symmetric spacecraft}
% source: cap11.tex l.178-384
\begin{qbox}
State the properties of the inertia matrix and give its form for spacecraft with a plane, an axis or a point of symmetry.
\qmeta{never asked (Ch.\,11: properties T\&N, only the statements).}
\end{qbox}
\begin{outline}
\begin{enumerate}
\item Symmetric (real eigenvalues, orthogonal eigenvectors), triangular inequalities, positive (semi)definite.
\item Principal axes: diagonalisation by a rotation.
\item Plane of symmetry: two zero products of inertia; axis: $\mathrm{diag}(I_1,I_1,I_3)$; point: $I_1\underline{\underline I}$.
\end{enumerate}
\end{outline}
\answer
\paragraph{Properties} of $\underline{\underline I}_C^{(B)}=\int\begin{bmatrix}y^2+z^2&-xy&-xz\\-yx&x^2+z^2&-yz\\-zx&-zy&x^2+y^2\end{bmatrix}dm$
(diagonal: moments of inertia; off-diagonal: products of inertia):
\begin{enumerate}
\item \textbf{symmetric}: real eigenvalues, real orthogonal eigenvectors; the eigenvectors form a rotation $\underline{\underline T}$ that diagonalises it:
$\underline{\underline T}^{-1}\underline{\underline I}\underline{\underline T}=\mathrm{diag}(\lambda_i)$ (\textbf{principal axes}, principal moments);
\item \textbf{triangular inequalities}: $I_{11}\le I_{22}+I_{33}$ and cyclic;
\item \textbf{positive (semi)definite}: $\underline v^T\underline{\underline I}\underline v=\int v^2r^2\sin^2\phi\,dm\ge0$.
\end{enumerate}
\paragraph{Symmetric spacecraft} ($\underline r$ from the centre of mass).
\begin{itemize}
\item \textbf{Plane of symmetry} $(x,z)$ (winged spacecraft): for each $dm$ at $(x,y,z)$ there is one at $(x,-y,z)$, so terms with $y$ cancel in pairs:
$\underline{\underline I}=\begin{bmatrix}I_{11}&0&I_{13}\\0&I_{22}&0\\I_{31}&0&I_{33}\end{bmatrix}$; the axis normal to the plane is principal.
\item \textbf{Axis of symmetry} $\hat b_3$ (wingless rocket): four symmetric masses cancel all products, $\underline{\underline I}=\mathrm{diag}(I_1,I_1,I_3)$;
any pair of transverse axes is principal.
\item \textbf{Point symmetry} (spherical satellite): eight symmetric masses: $I_1\underline{\underline I}_{3\times3}$; any axes are principal.
\end{itemize}
\begin{keyf}
$\underline{\underline I}^T=\underline{\underline I}$, $I_{11}\le I_{22}+I_{33}$, $\underline v^T\underline{\underline I}\underline v\ge0$; \ plane $(x,z)$: $I_{12}=I_{23}=0$; \ axis: $\mathrm{diag}(I_1,I_1,I_3)$; \ sphere: $I_1\underline{\underline I}$.
\end{keyf}
\pitfalls
\begin{itemize}
\item Two equal moments do not require axial symmetry (e.g.\ a square-section prism).
\end{itemize}
\source{cap11.tex l.178--227 (properties), l.276--287 (principal axes), l.329--384 (symmetric spacecraft). Formulario: \fml{11-IPR}, \fml{11-SYM}.}

% ---------------------------------------------------------------------
\question{E18}{Torque-free motion of an axisymmetric body}
% source: cap11.tex l.475-678
\begin{qbox}
Describe the torque-free motion of an axisymmetric rigid body: angular velocity in the body frame, Euler angles, oblate and prolate bodies.
\qmeta{never asked (Ch.\,11, T\&N).}
\end{qbox}
\begin{outline}
\begin{enumerate}
\item $I_1=I_2=I_T$: $\omega_3$ const, $\omega_{1,2}$ harmonic with $\omega_p=\omega_{30}(I_3/I_T-1)$: $\underline\omega$ describes a cone around $\hat e_3$.
\item $\hat E_3\parallel\underline H_C$: $\theta$ constant, $\tan\theta=I_T\omega_{12}/(I_3\omega_{30})$, $\dot\psi=I_T\omega_{12}^2/(H_Cs_\theta^2)>0$, $\dot\phi=-\omega_p$.
\item Oblate ($I_3>I_T$): retrograde ($\dot\phi<0$); prolate: prograde. Body and space cones.
\end{enumerate}
\end{outline}
\answer
\paragraph{Angular velocity in body axes.} With $I_1=I_2=I_T\neq I_3$ and no torque: $I_3\dot\omega_3=0$ $\Rightarrow$ $\omega_3=\omega_{30}$, and with
$\omega_p:=\omega_{30}\left(\frac{I_3}{I_T}-1\right)$: $\dot\omega_1=-\omega_p\omega_2$, $\dot\omega_2=\omega_p\omega_1$ $\Rightarrow$ $\ddot\omega_{1,2}+\omega_p^2\omega_{1,2}=0$:
\begin{equation}
\omega_1=\omega_{12}\cos(\omega_pt+\varphi),\quad \omega_2=\omega_{12}\sin(\omega_pt+\varphi),\quad \omega_3=\omega_{30},\qquad \omega_{12}=\sqrt{\omega_{10}^2+\omega_{20}^2}.
\end{equation}
In the body frame $\underline\omega$ describes a cone about $\hat e_3$ at rate $\omega_p$ (counter-clockwise if $\omega_p>0$).
\paragraph{Inertial description.} Choose $\hat E_3\parallel\underline H_C$ (constant) and Euler angles 3-1-3 with $\omega_{30}\ge0$. The last column of $\underline{\underline R}$ gives
$c_\theta=I_3\omega_3/H_C$ $\Rightarrow$ \textbf{$\theta$ constant} (nutation angle between $\hat e_3$ and $\underline H_C$), $\tan\theta=\frac{I_T\omega_{12}}{I_3\omega_{30}}$; and
\begin{equation}
\dot\psi=\frac{I_T\omega_{12}^2}{H_Cs_\theta^2}>0,\qquad \dot\theta=0,\qquad \dot\phi=-\omega_p=-\omega_{30}\left(\frac{I_3}{I_T}-1\right).
\end{equation}
$\underline\omega=\dot\psi\hat E_3+\dot\phi\hat e_3$: $\hat e_3$ and $\underline\omega$ rotate counter-clockwise about $\underline H_C$ (precession). \textbf{Prolate} body ($I_3<I_T$): $\dot\phi>0$,
prograde precession, the body cone rolls on the outside of the space cone. \textbf{Oblate} body ($I_3>I_T$): $\dot\phi<0$, retrograde precession, the body cone
contains the space cone and rolls on it with its inner side. With dissipation (\qref{D9}) oblate bodies tend to pure spin, prolate bodies to flat spin.
\begin{keyf}
$\omega_p=\omega_{30}(I_3/I_T-1)$; \ $\tan\theta=\dfrac{I_T\omega_{12}}{I_3\omega_{30}}$; \ $\dot\psi=\dfrac{I_T\omega_{12}^2}{H_Cs_\theta^2}$; \ $\dot\phi=-\omega_p$.
\end{keyf}
\pitfalls
\begin{itemize}
\item In the body frame the rate is $\omega_p$; in inertial space the precession rate is $\dot\psi$ (different).
\end{itemize}
\source{cap11.tex l.475--678 (axisymmetric torque-free motion), l.680--707 (general body). Formulario: \fml{11-AXS}, \fml{11-ANG}.}

% ---------------------------------------------------------------------
\question{E19}{Attitude manoeuvres of spinning satellites}
% source: cap11.tex l.979-1100
\begin{qbox}
Describe spin-up/spin-down manoeuvres (thrusters and yo-yo despin) and impulsive attitude manoeuvres of spinning satellites.
\qmeta{never asked (Ch.\,11: T; spin-up/down T\&N).}
\end{qbox}
\begin{outline}
\begin{enumerate}
\item Spin about the max-inertia axis: intrinsic stabilisation; manoeuvres to change spin rate or orientation.
\item Thruster pairs: $\underline L_C=2RF\hat e_3$, no translation; $\omega_3=\omega_{30}\pm\frac{2RF}{I_3}t$.
\item Yo-yo: $H$ conserved, $\omega_{3f}=\frac{I_3\omega_{30}}{I_3'+2ml^2}$.
\item Impulsive torque: $\Delta\omega_3=\frac{L_3}{I_3}\Delta t$, attitude unchanged.
\end{enumerate}
\end{outline}
\answer
\paragraph{Motivation.} A satellite spinning about its maximum-inertia axis is intrinsically stabilised; manoeuvres are needed to spin it up or
down and to change its orientation.
\paragraph{Spin-up/spin-down with thrusters.} Thrusters are fired \textbf{in pairs} with opposite forces at distance $R$ from the axis: no effect on
the centre of mass, torque $\underline L_C=\underline p_1\times\underline F_1+\underline p_2\times\underline F_2=2RF\hat e_3$. With $\omega_1=\omega_2=0$:
$I_3\dot\omega_3=\pm2RF$, $\omega_3=\omega_{30}\pm\frac{2RF}{I_3}t$ (constant $F$); the attitude follows from the kinematic equations.
\paragraph{Yo-yo despin.} Two equal masses $m$ on cables of length $l$ are deployed symmetrically (centre of mass unchanged); no external torque, so
$H$ is conserved: $I_3\omega_{30}=(I_3'+2ml^2)\omega_{3f}$,
\begin{equation}
\omega_{3f}=\frac{I_3\,\omega_{30}}{I_3'+2ml^2}<\omega_{30},
\end{equation}
($I_3'$ = moment of inertia without the masses), more effective for longer cables.
\paragraph{Impulsive attitude manoeuvres.} For a spin-up with constant $L_3$: $\omega_{3f}=\omega_{30}+\frac{L_3}{I_3}\Delta t$,
$\phi_f=\phi_0+\omega_{30}\Delta t+\frac12\frac{L_3}{I_3}\Delta t^2$. The impulsive limit $\Delta t\to0$ with $\frac{L_3}{I_3}\Delta t=\Delta\omega_3$ finite (torque as a
Dirac delta) gives: (a) an instantaneous change of angular velocity $\Delta\omega_3$; (b) no change of attitude ($\phi_f\to\phi_0$) -- the rotational analogue of
the impulsive-thrust approximation (\qref{A5}). Reorientation of the spin axis by $\delta$ uses two such transverse impulses separated by half a
precession period.
\begin{keyf}
$L_C=2RF$, $\omega_3=\omega_{30}\pm\frac{2RF}{I_3}t$; \ yo-yo $\omega_{3f}=\dfrac{I_3\omega_{30}}{I_3'+2ml^2}$; \ impulsive: $\Delta\omega=\dfrac{L}{I}\Delta t$, attitude unchanged.
\end{keyf}
\pitfalls
\begin{itemize}
\item The yo-yo works by angular-momentum conservation (no propellant).
\end{itemize}
\source{cap11.tex l.979--1054 (spin-up/down, yo-yo), l.1056--1096 (impulsive manoeuvres), l.1098--1195 (reorientation, N). Formulario: \fml{11-SPN}, \fml{11-IMP}.}
